\newpage

\section{Стохастический блочно-координатный Frank--Wolfe}
\label{sec:scfw}

В этой главе рассматривается стохастический вариант Frank--Wolfe для крупномасштабной задачи равновесного распределения потоков~\cite{ignashin2026stochastic}.
Этот вариант удобно называть Stochastic Origin Frank--Wolfe (SOFW).
В контексте данной работы он рассматривается как специальный случай блочно-координатного FW (BCFW) со стохастическим выбором блоков.
Его ключевая идея состоит в выборе случайных блоков и обновлении только части координат на каждой итерации.

Классический FW в транспортной задаче требует на каждой итерации найти кратчайшие пути от всех источников ко всем стокам.
Если граф содержит много зон и узлов, стоимость одной итерации становится основной вычислительной проблемой.
Стохастический подход уменьшает стоимость итерации: вместо полного линейного оракула выбирается случайное подмножество источников или OD-пар, а остальные блоки фиксируются.

\subsection{Блочная структура транспортной задачи}

Рассмотрим вектор потоков по путям как конкатенацию блоков по OD-парам:
\[
    x=
    \left(x^{(i,j)}\right)_{(i,j)\in OD},
    \qquad
    x^{(i,j)}\in\R^{|\cP_{i,j}|}.
\]
Каждый блок соответствует множеству путей $\cP_{i,j}$ между источником $i$ и стоком $j$.
Для блока выполняются ограничения спроса
\[
    \sum_{p\in\cP_{i,j}}x_p=d_{ij},
    \qquad
    x_p\geq 0.
\]
Следовательно, допустимое множество потоков по путям является декартовым произведением:
\[
    \cX=
    \prod_{(i,j)\in OD}\cX^{(i,j)}.
\]
Переход от потоков по путям к потокам по ребрам задается линейным отображением
\[
    F(x)=\Theta x,
    \qquad
    f=F(x).
\]

В классическом FW линейный оракул решается по всем блокам:
\[
    s_k\in
    \argmin_{x\in\cX}
    \inner{\nabla\Psi(F(x_k))}{F(x)}.
\]
Так как
\[
    \inner{\nabla\Psi(F(x_k))}{F(x)}
    =
    \sum_{(i,j)\in OD}
    \inner{\tau(f_k)}{\Theta_{ij}x^{(i,j)}},
\]
оракул распадается на независимые подзадачи по OD-парам.
Каждая такая подзадача является назначением всего спроса на кратчайший путь при текущих дорожных временах.

\subsection{Stochastic Origin Frank--Wolfe}

На итерации $k$ SOFW выбирает случайное подмножество OD-пар $\cI_k\subset OD$ и обновляет только соответствующие блоки.
Остальные блоки фиксируются в текущих значениях:
\[
    \cX_k
    =
    \left\{
        x\in\cX
        \;\middle|\;
        x^{(i,j)}=x_k^{(i,j)}
        \text{ для всех }(i,j)\notin\cI_k
    \right\}.
\]
Линейный оракул решается на суженном множестве:
\[
    s_k
    \in
    \argmin_{x\in\cX_k}
    \inner{\nabla\Psi(f_k)}{F(x)}.
\]
На практике удобно выбирать не отдельные OD-пары, а подмножество источников.
Если выбран источник $i$, то в батч включаются все корреспонденции $(i,j)$.
Это позволяет один раз запустить алгоритм Дейкстры от источника $i$ и получить кратчайшие пути до всех стоков.

\subsection{Статическое назначение потоков и роль модели Бекмана}

SOFW мотивируется задачей прогнозирования транспортных потоков.
В транспортном планировании требуется оценить распределение движения по дорожной сети
при заданной матрице спроса.
В статической модели считается, что спрос за рассматриваемый период фиксирован,
а равновесие описывается принципом Вардропа: все используемые пути между одной OD-парой
имеют минимальное время движения, а неиспользуемые пути не короче используемых.

Модель Бекмана дает выпуклую оптимизационную формулировку этого равновесия:
\[
    \min_{f=\Theta x,\;x\in\cX}
    \Psi(f)=
    \sum_{e\in\cE}\int_0^{f_e}\tau_e(z)\,dz.
\]
Если функции $\tau_e$ неубывающие, то задача выпукла.
Оптимальность $x^\ast$ означает, что для любой допустимой перераспределенной траектории
$x$ выполнено вариационное неравенство
\[
    \inner{\tau(f^\ast)}{\Theta x-\Theta x^\ast}\geq 0.
\]
Это ровно условие равновесия: нельзя уменьшить суммарную линейную стоимость при текущих
временах движения.

Классический FW применим здесь особенно естественно.
Линейная подзадача
\[
    \min_{x\in\cX}\inner{\tau(f^k)}{\Theta x}
\]
распадается по OD-парам:
\[
    \min_{x^{(i,j)}\in\cX^{(i,j)}}
    \sum_{p\in\cP_{i,j}}x_p
    \sum_{e\in p}\tau_e(f^k_e).
\]
Оптимум достигается назначением всего спроса $d_{ij}$ на кратчайший путь.
Поэтому каждая полная итерация FW фактически состоит из набора запусков алгоритма
кратчайших путей.

\subsection{Path-based и link-based представления}

В теории удобно работать с путевыми потоками $x$, потому что допустимое множество имеет
явную блочную структуру:
\[
    \cX=\prod_{(i,j)\in OD}
    \left\{
        x^{(i,j)}\geq 0\;\middle|\;
        \sum_{p\in\cP_{i,j}}x_p=d_{ij}
    \right\}.
\]
Однако число возможных путей может быть экспоненциальным.
Поэтому в реализации обычно не хранится весь вектор $x$.
Хранятся потоки по ребрам $f$, а очередное решение линейного оракула строится как
all-or-nothing assignment: для каждой OD-пары весь спрос временно отправляется на текущий
кратчайший путь.

SOFW требует больше информации, чем полный FW.
Если обновляется только часть источников, нужно уметь вычесть старый вклад этих источников
из суммарного потока и добавить новый.
Поэтому алгоритм хранит разложение
\[
    f=\sum_{o\in O}f^{(o)},
\]
где $f^{(o)}$ -- поток по ребрам, порожденный спросом из источника $o$.
Когда выбран батч источников $\mathcal A_k$, обновляются только компоненты
$f^{(o)}$, $o\in\mathcal A_k$, а остальные остаются фиксированными.
Именно это разложение объясняет главный компромисс метода: SOFW уменьшает время итерации,
но требует больше памяти.

\begin{algorithm}
\caption{Stochastic Origin Frank--Wolfe}
\label{alg:sofw}
\begin{algorithmic}
\REQUIRE допустимое множество $\cX=\prod_{w\in OD}\cX^w$, отображение $F(x)=\Theta x$, вероятности выбора блоков, шаги $\gamma_k$
\STATE выбрать начальное $x_0\in\cX$ и положить $f_0=F(x_0)$
\FOR{$k=0,1,2,\ldots$}
\STATE выбрать случайное подмножество $\cI_k\subset OD$
\STATE определить суженное множество
\[
    \cX_k=
    \left(\prod_{w\in\cI_k}\cX^w\right)
    \times
    \left(\prod_{w\notin\cI_k}\{x_k^w\}\right)
\]
\STATE вычислить $s_k\in\argmin_{x\in\cX_k}\inner{\nabla\Psi(f_k)}{F(x)}$
\STATE вычислить направление $d_k=F(s_k)-F(x_k)$ или его взвешенную несмещенную оценку
\STATE обновить $f_{k+1}=f_k+\gamma_kd_k$
\STATE обновить $x_{k+1}=x_k+\gamma_k(s_k-x_k)$
\ENDFOR
\end{algorithmic}
\end{algorithm}

\subsection{Взвешенный вариант}

Во взвешенном варианте источники выбираются с вероятностями, пропорциональными суммарному спросу.
Пусть
\[
    d_i=\sum_j d_{ij},
    \qquad
    \cC_i=\{(i,j)\in OD\}.
\]
Здесь $d_i$ -- полный спрос, исходящий из источника $i$, а $\cC_i$ -- множество всех
корреспонденций с этим источником.
В качестве веса источника можно взять нормированный спрос
\[
    w_i
    =
    \frac{d_i}{\sum_{\ell}d_\ell},
    \qquad
    \sum_i w_i=1.
\]
Тогда источник $i$ выбирается с вероятностью $w_i$.
Если после выбора источника корреспонденция внутри $\cC_i$ выбирается равномерно, то
вес OD-блока $(i,j)$ равен
\[
    w_{ij}
    =
    \mathbb{P}\{(i,j)\}
    =
    w_i
    \cdot
    \frac{1}{|\cC_i|}
    =
    \frac{d_i}{\sum_{\ell}d_\ell}
    \cdot
    \frac{1}{|\cC_i|}.
\]
Такой выбор чаще обновляет источники с большим суммарным спросом.
На практике можно выбирать сразу батч источников с вероятностями $w_i$ и затем включать
в батч все корреспонденции выбранных источников; это соответствует той же идее, но лучше
согласуется с вычислением кратчайших путей от одного источника ко всем стокам.

Для сохранения несмещенности направления используется importance-weighted коррекция:
\[
    d_k
    =
    \frac{1}{|\cI_k|}
    W_k
    \sum_{(i,j)\in\cI_k}
    \frac{1}{w_{ij}}
    \left(
        F(s_k^{[(i,j)]})-F(x_k^{[(i,j)]})
    \right),
    \qquad
    W_k=
    \left(
        \sum_{(i,j)\in\cI_k}\frac{1}{w_{ij}}
    \right)^{-1}.
\]

\subsection{Сходимость}

Сходимость SOFW следует из анализа Block-Coordinate Frank--Wolfe на произведениях множеств~\cite{lacoste2013block,jaggi2012blockcoordfw}.
Для каждого блока $w=(i,j)$ вводится блочная константа кривизны:
\[
C_\Psi^w
=
\sup
\frac{2}{\gamma^2}
\left[
    \Psi(F(y))
    -
    \Psi(F(x))
    -
    \inner{\nabla\Psi(F(x))}{F(y^{[w]})-F(x^{[w]})}
\right],
\]
где супремум берется по $x\in\cX$, $s^w\in\cX^w$, $\gamma\in[0,1]$ и
\[
    y=x+\gamma(s^{[w]}-x^{[w]}).
\]
Суммарная блочная кривизна равна
\[
    C_\Psi^\otimes=\sum_{w\in OD}C_\Psi^w.
\]

Если на каждой итерации выбирается одна OD-пара равномерно случайно, то для ожидаемой субоптимальности выполняется оценка
\begin{equation}
    \E[\Psi(f_k)]-\Psi(f^\ast)
    \leq
    \frac{2|OD|}{k+2|OD|}
    \left[
        C_\Psi^\otimes+\Psi(f_0)-\Psi(f^\ast)
    \right].
    \label{eq:sofw-rate-single}
\end{equation}
Для батчевого варианта, где на итерации выбирается подмножество $\cI_k$ размера $m$, направление можно записать как
\[
    d_k=
    \frac{1}{m}
    \sum_{w\in\cI_k}
    \left(s_k^{[w]}-x_k^{[w]}\right),
    \qquad
    x_{k+1}=x_k+\gamma_kd_k.
\]
При равномерном выборе без возвращения условное математическое ожидание прогресса имеет тот же вид, что и в небатчевом случае:
\[
    \E[\Psi(F(x_{k+1}))\mid x_k]
    \leq
    \Psi(F(x_k))
    -
    \frac{\gamma_k}{|OD|}g(x_k)
    +
    \frac{\gamma_k^2}{2|OD|}C_\Psi^\otimes.
\]

\begin{theorem}
\label{thm:sofw-batched-rate}
Пусть $\{x_k\}$ порождается батчевым SOFW с произвольным размером батча $m$ и шагами
\[
    \gamma_k=\frac{2|OD|}{k+2|OD|}.
\]
Тогда
\[
    \E[\Psi(F(x_k))]-\Psi(F(x^\ast))
    \leq
    \frac{2|OD|}{k+2|OD|}
    \left[
        C_\Psi^\otimes+\Psi(F(x_0))-\Psi(F(x^\ast))
    \right].
\]
Следовательно, для достижения $\varepsilon$-точности в ожидании достаточно
\[
    \mathcal{O}\left(\frac{|OD|}{\varepsilon}\right)
\]
итераций.
\end{theorem}

Батчирование в такой нормировке не улучшает теоретический порядок по числу итераций, но уменьшает дисперсию направления и может существенно улучшать время работы.
В транспортной задаче это особенно важно, поскольку стоимость итерации почти линейно зависит от числа источников, для которых запускается поиск кратчайших путей.

\subsection{Доказательство батчевой оценки}

Доказательство для батчевого SOFW начинается с явного разложения полного FW-зазора по
OD-блокам.
Пусть выбран блок $w\in OD$ и
\[
    s^w(x)\in
    \argmin_{u\in\cX^w}
    \inner{\nabla\Psi(F(x))}{F(u^{[w]})}.
\]
Здесь $u^{[w]}$ означает вектор, у которого активен только блок $w$.
Определим блочный FW-зазор
\[
    g^w(x)
    =
    \inner{\nabla\Psi(F(x))}
    {F\left((x^w-s^w(x))^{[w]}\right)}
    \geq 0.
\]
Тогда полный FW-зазор раскладывается в сумму
\[
    g(x)
    =
    \sum_{w\in OD}g^w(x).
\]
Это равенство является ключевым: линейный оракул в модели Бекмана распадается по
OD-парам, поэтому потеря от обновления только части блоков выражается через долю
неиспользованных слагаемых в этой сумме.

Обновление только этого блока имеет вид
\[
    x^+=x+\gamma(s^w(x)-x^w)^{[w]}.
\]
По определению блочной кривизны
\[
    \Psi(F(x^+))
    \leq
    \Psi(F(x))
    -
    \gamma g^w(x)
    +
    \frac{\gamma^2}{2}C_\Psi^w.
\]
Если блок выбирается равномерно, то после взятия условного ожидания получается
\[
    \E[g^w(x)\mid x]
    =
    \frac{1}{|OD|}\sum_{w\in OD}g^w(x)
    =
    \frac{1}{|OD|}g(x).
\]
Аналогично
\[
    \E[C_\Psi^w]=\frac{1}{|OD|}C_\Psi^\otimes.
\]
Отсюда следует рекурсия
\[
    \E[\Psi(F(x_{k+1}))\mid x_k]
    \leq
    \Psi(F(x_k))
    -
    \frac{\gamma_k}{|OD|}g(x_k)
    +
    \frac{\gamma_k^2}{2|OD|}C_\Psi^\otimes.
\]

Для батча размера $m$ направление записывается как среднее по выбранным блокам:
\[
    d_k
    =
    \frac{1}{m}
    \sum_{w\in\cI_k}
    (s^w(x_k)-x_k^w)^{[w]}.
\]
Поскольку $F$ линейно, а $\Psi$ выпукла, точку после батчевого шага можно сравнить со
средним одноблочных шагов:
\[
    x_k+\gamma d_k
    =
    \frac{1}{m}
    \sum_{w\in\cI_k}
    \left(
        x_k+\gamma(s^w(x_k)-x_k^w)^{[w]}
    \right).
\]
Следовательно,
\[
    \Psi(F(x_k+\gamma d_k))
    \leq
    \frac{1}{m}
    \sum_{w\in\cI_k}
    \Psi\left(
        F\left(x_k+\gamma(s^w(x_k)-x_k^w)^{[w]}\right)
    \right).
\]
Подставляя одноблочную оценку в правую часть, получаем
\[
    \Psi(F(x_k+\gamma d_k))
    \leq
    \Psi(F(x_k))
    -
    \frac{\gamma}{m}
    \sum_{w\in\cI_k}g^w(x_k)
    +
    \frac{\gamma^2}{2m}
    \sum_{w\in\cI_k}C_\Psi^w.
\]
Если блоки выбираются равномерно без возвращения, то каждый блок попадает в батч с
вероятностью $m/|OD|$.
Поэтому
\[
    \E\left[
        \frac{1}{m}\sum_{w\in\cI_k}g^w(x_k)
        \;\middle|\;x_k
    \right]
    =
    \frac{1}{|OD|}\sum_{w\in OD}g^w(x_k)
    =
    \frac{1}{|OD|}g(x_k),
\]
и аналогично
\[
    \E\left[
        \frac{1}{m}\sum_{w\in\cI_k}C_\Psi^w
        \;\middle|\;x_k
    \right]
    =
    \frac{1}{|OD|}C_\Psi^\otimes.
\]
Именно поэтому теоретическая рекурсия для батча совпадает с одноблочной рекурсией.
Размер батча не входит в порядок сходимости, но влияет на дисперсию направления и на
реальное время достижения заданного зазора.

Переход от рекурсии к явной оценке использует выпуклость:
\[
    g(x_k)\geq \Psi(F(x_k))-\Psi(F(x^\ast)).
\]
Обозначив
\[
    h_k=\E[\Psi(F(x_k))-\Psi(F(x^\ast))],
\]
получаем
\[
    h_{k+1}
    \leq
    \left(1-\frac{\gamma_k}{|OD|}\right)h_k
    +
    \frac{\gamma_k^2}{2|OD|}C_\Psi^\otimes.
\]
При выборе
\[
    \gamma_k=\frac{2|OD|}{k+2|OD|}
\]
индукция дает
\[
    h_k
    \leq
    \frac{2|OD|}{k+2|OD|}
    \left[
        C_\Psi^\otimes+h_0
    \right].
\]
Именно эта формула и записана в теореме~\ref{thm:sofw-batched-rate}.

\subsection{Вычислительная сложность и память}

Пусть $n$ есть число источников, $V=|\cV|$ число узлов, а $E=|\cE|$ число ребер графа.
Полная итерация FW требует запустить поиск кратчайших путей из всех источников и имеет оценку
\[
    \mathcal{O}\left(n(E+V\log V)\right).
\]
Если SOFW выбирает долю $\alpha$ источников, то стоимость итерации становится
\[
    \mathcal{O}\left(\alpha n(E+V\log V)\right).
\]
Цена такого ускорения состоит в увеличенном потреблении памяти: необходимо хранить разложение суммарного потока по источникам
\[
    f=\sum_{a=1}^{n}f^{(a)}.
\]
Если за итерацию выбирается множество источников $\mathcal{A}$, то увеличение памяти оценивается примерно как
\[
    \frac{n}{|\mathcal{A}|}.
\]
Таким образом, SOFW явно задает компромисс между временем итерации, качеством стохастического направления и расходом памяти.

\subsection{Экспериментальные результаты}

Эксперименты проводились на транспортных сетях из открытого репозитория~\cite{bstabler}.
Сравнивались классический FW, равномерный SOFW и взвешенный SOFW-w.
Для крупных графов стохастический метод давал существенно меньший зазор Frank--Wolfe при том же вычислительном бюджете.

\begin{table}[H]
\centering
\caption{Характеристики транспортных датасетов и оценка стоимости полной итерации $\widehat{\mathcal{C}}_{\mathrm{iter}}=Z(E+N\ln N)$.}
\label{tab:sofw-datasets}
\begin{tabular}{lcccc}
\toprule
Датасет & Зоны $Z$ & Ребра $E$ & Узлы $N$ & $\widehat{\mathcal{C}}_{\mathrm{iter}}$ \\
\midrule
SiouxFalls & 24 & 76 & 24 & $3.66\cdot 10^3$ \\
Anaheim & 38 & 914 & 416 & $1.30\cdot 10^5$ \\
Terrassa-Asymmetric & 55 & 3264 & 1609 & $8.33\cdot 10^5$ \\
Winnipeg & 147 & 2836 & 1052 & $1.49\cdot 10^6$ \\
Chicago-Sketch & 387 & 2950 & 933 & $3.61\cdot 10^6$ \\
GoldCoast & 1068 & 11140 & 4807 & $5.54\cdot 10^7$ \\
Birmingham-England & 898 & 33937 & 14639 & $1.57\cdot 10^8$ \\
Chicago-Regional & 1790 & 39018 & 12982 & $2.90\cdot 10^8$ \\
Philadelphia & 1525 & 40003 & 13389 & $2.55\cdot 10^8$ \\
\bottomrule
\end{tabular}
\end{table}

\begin{figure}[H]
    \centering
    \subfloat[Chicago-Sketch]{\includegraphics[width=0.48\textwidth]{scfw/Chicago-Sketch.pdf}}
    \subfloat[Birmingham-England]{\includegraphics[width=0.48\textwidth]{scfw/Birmingham-England.pdf}}
    \caption{Сравнение SOFW и FW на средних транспортных сетях.}
    \label{fig:sofw-medium-1}
\end{figure}

\begin{figure}[H]
    \centering
    \subfloat[Chicago-Regional]{\includegraphics[width=0.48\textwidth]{scfw/chicago-regional.pdf}}
    \subfloat[GoldCoast]{\includegraphics[width=0.48\textwidth]{scfw/GoldCoast.pdf}}
    \caption{SOFW на крупных сетях: выигрыш увеличивается вместе с размером графа.}
    \label{fig:sofw-medium-2}
\end{figure}

\begin{figure}[H]
    \centering
    \includegraphics[width=0.8\textwidth]{scfw/Philadelphia.pdf}
    \caption{Датасет Philadelphia: сравнение SOFW и FW на крупной транспортной сети.}
    \label{fig:sofw-philadelphia}
\end{figure}

\begin{figure}[H]
    \centering
    \subfloat[SiouxFalls]{\includegraphics[width=0.48\textwidth]{scfw/SiouxFalls.pdf}}
    \subfloat[Anaheim]{\includegraphics[width=0.48\textwidth]{scfw/Anaheim.pdf}}
    \caption{SOFW и FW на малых сетях: преимущество стохастического подхода менее выражено из-за накладных расходов.}
    \label{fig:sofw-small}
\end{figure}

\begin{figure}[H]
    \centering
    \subfloat[Terrassa-Asymmetric]{\includegraphics[width=0.48\textwidth]{scfw/Terrassa-Asymmetric.pdf}}
    \subfloat[Winnipeg]{\includegraphics[width=0.48\textwidth]{scfw/Winnipeg.pdf}}
    \caption{Переходная зона между малыми и средними сетями: стохастический подход начинает давать устойчивый выигрыш.}
    \label{fig:sofw-small-2}
\end{figure}

\subsection{Интерпретация графиков SOFW}

На малых сетях SiouxFalls и Anaheim стоимость полного запуска кратчайших путей невелика.
Поэтому экономия от частичного оракула не всегда перекрывает дополнительные расходы:
нужно хранить разложение потока по источникам и аккуратно обновлять только выбранные
компоненты.
На таких задачах классический FW остается сильным базовым методом.

На Terrassa-Asymmetric и Winnipeg картина меняется.
Число зон и ребер уже достаточно велико, чтобы стоимость полной итерации стала заметной.
SOFW начинает выигрывать по времени, особенно на начальном участке, где даже грубое
стохастическое направление быстро снижает зазор.

На Chicago-Sketch, GoldCoast, Birmingham-England, Chicago-Regional и Philadelphia преимущество
становится основным практическим результатом.
Чем больше сеть, тем менее рационально на каждой итерации решать линейный оракул по всем
источникам.
Случайное обновление части источников дает больше итераций за тот же бюджет времени.
Даже если каждая отдельная итерация менее точна, суммарный прогресс оказывается лучше.

Взвешенный вариант SOFW-w особенно полезен на ранних итерациях.
Если источник создает большой суммарный спрос, его вклад в поток по ребрам велик.
Обновление такого источника сильнее меняет функцию Бекмана и быстрее уменьшает грубую
начальную ошибку.
По мере приближения к равновесию различие между равномерным и взвешенным выбором уменьшается:
локальные корректировки становятся тоньше, а роль случайности и конкретных OD-пар возрастает.

Основные экспериментальные выводы:
\begin{enumerate}
    \item SOFW существенно превосходит FW на крупных графах.
    \item Преимущество SOFW растет вместе с размером сети.
    \item На малых сетях метод остается конкурентоспособным, но накладные расходы хранения разложения потока уменьшают выигрыш.
    \item Взвешенный SOFW-w быстрее уменьшает зазор на первых итерациях.
    \item Долгосрочно равномерный и взвешенный варианты дают близкое качество, особенно на крупных графах.
\end{enumerate}

\subsection{Переход от OD-пар к источникам}

Теоретически блоком может быть отдельная OD-пара $w=(i,j)$.
Однако в дорожном графе эффективнее выбирать источники.
Если выбран источник $i$, то один запуск алгоритма кратчайших путей из $i$ дает расстояния
до всех стоков $j$.
Поэтому батч источников
\[
    \mathcal A_k\subset O
\]
порождает батч OD-пар
\[
    \cI_k=\{(i,j)\in OD\mid i\in\mathcal A_k\}.
\]
Такой выбор лучше соответствует структуре вычислений:
стоимость одной итерации примерно пропорциональна числу выбранных источников, а не числу
отдельных OD-пар.

\subsection{Обновление разложения потока}

Пусть
\[
    f_k=\sum_{i\in O}f_k^{(i)}.
\]
Для выбранного источника $i$ строится новый all-or-nothing поток $\hat f_k^{(i)}$:
весь спрос из $i$ назначается на кратчайшие пути при текущих временах $\tau(f_k)$.
Затем направление для этого источника равно
\[
    d_k^{(i)}=\hat f_k^{(i)}-f_k^{(i)}.
\]
Если выбран батч $\mathcal A_k$, суммарное направление записывается как
\[
    d_k=\sum_{i\in\mathcal A_k}d_k^{(i)}.
\]
Обновление выполняется только для выбранных источников:
\[
    f_{k+1}^{(i)}
    =
    \begin{cases}
        f_k^{(i)}+\gamma_k d_k^{(i)}, & i\in\mathcal A_k,\\
        f_k^{(i)}, & i\notin\mathcal A_k.
    \end{cases}
\]
После этого
\[
    f_{k+1}=\sum_{i\in O}f_{k+1}^{(i)}.
\]
Эта формула показывает, почему SOFW требует хранения $f^{(i)}$.
Без такого разложения нельзя локально заменить вклад выбранного источника, не пересчитывая
полный поток.

\subsection{Равномерная и взвешенная выборка}

При равномерной выборке каждый источник имеет одинаковую вероятность попасть в батч.
Это просто и дает чистую теоретическую оценку.
Но спрос по источникам может быть сильно неравномерным.
Если источник генерирует большой поток, его обновление может сильнее уменьшить функционал.
Поэтому рассматривается взвешенная версия.

Пусть
\[
    D_i=\sum_j d_{ij}
\]
-- полный спрос из источника $i$.
Вероятность выбора источника задается как
\[
    p_i=\frac{D_i}{\sum_{\ell\in O}D_\ell}.
\]
Чтобы направление не стало смещенным в сторону крупных источников, можно использовать
importance weighting.
В простейшей записи вклад выбранного источника корректируется множителем $1/p_i$:
\[
    \tilde d_k^{(i)}=\frac{1}{p_i}d_k^{(i)}.
\]
Тогда
\[
    \E[\tilde d_k^{(i)}]=\sum_{i\in O}d_k^{(i)}
\]
с точностью до нормировки батча.
На практике эта коррекция стабилизирует сравнение равномерного и взвешенного SOFW.

\subsection{Влияние батча на теоретический порядок}

Интуитивно может показаться, что батч размера $m$ должен ускорять сходимость в $m$ раз.
Однако в нормированной записи направление усредняется по выбранным блокам.
Ожидаемый прогресс на одну итерацию остается пропорционален полному FW-зазору, деленному
на число блоков:
\[
    \E[\Delta_k\mid x_k]\approx -\frac{\gamma_k}{|OD|}g(x_k).
\]
Батч уменьшает дисперсию, но не меняет этот множитель в базовой оценке.
Практический выигрыш появляется через время выполнения:
если батч можно обработать значительно быстрее полного оракула, то за тот же бюджет времени
алгоритм делает больше обновлений.

\subsection{Память и время}

Если в графе $|O|$ источников и $|\cE|$ ребер, хранение полного потока требует
\[
    \cO(|\cE|)
\]
памяти.
Хранение разложения по источникам требует
\[
    \cO(|O||\cE|)
\]
в худшем случае.
На практике потоки могут быть разреженными, но порядок показывает потенциальную цену SOFW.
Время одной полной итерации FW можно оценить как
\[
    |O|\cdot T_{\mathrm{SP}},
\]
где $T_{\mathrm{SP}}$ -- стоимость одного запуска кратчайших путей.
Если SOFW выбирает долю $\alpha$ источников, то итерация стоит примерно
\[
    \alpha |O|\cdot T_{\mathrm{SP}},
\]
но требует поддерживать разложение $f=\sum_i f^{(i)}$.

\subsection{Итоги раздела}

Стохастический блочно-координатный подход показывает, что для крупных транспортных сетей главная стоимость FW -- это
не формула шага и не вычисление градиента, а полный линейный оракул.
SOFW уменьшает эту стоимость, выбирая только часть источников.
Теоретическая оценка остается такой же по порядку, как у блочного FW, а практические графики
показывают, что на больших сетях экономия времени важнее шума стохастического направления.


\subsection{Выводы по BCFW/SOFW}

Стохастический блочно-координатный подход переносит идею SGD на Frank--Wolfe: вместо полного дорогого оракула используется случайный частичный оракул.
В транспортной задаче такой выбор особенно естественен, поскольку линейная подзадача распадается по источникам и OD-парам.
BCFW/SOFW позволяет масштабировать проекционно-свободную оптимизацию на крупные сети, но требует аккуратного учета компромисса между памятью, числом обновляемых блоков и качеством направления.
Эта же блочная идея используется в следующей главе для построения децентрализованного алгоритма.
